\section{自监督学习回顾}

本节课的开头简要回顾了上一节课讨论的自监督学习（Self-Supervised Learning）内容。自监督学习的核心思想是：在没有人工标签的情况下，通过设计某种 pretext task，让模型从大量无标签数据中学习到有用的特征表示。

\begin{figure}[H]
\centering
\includegraphics[width=0.95\textwidth]{slides-images/slide-001.jpg}
\caption{自监督学习的基本范式：先在大规模无标签数据上通过 pretext task 训练 encoder-decoder，然后将 encoder 迁移到下游有标签任务\protect\footnotemark}
\end{figure}
\footnotetext{Slides 第2页。}

\subsection{对比学习（Contrastive Learning）}

对比学习是自监督学习中非常成功的一个范式。其核心思想是：

\begin{itemize}
    \item 对每张输入图片施加两种随机增强（Random Augmentation），得到一对正样本
    \item 来自不同图片的增强视为负样本
    \item 训练网络拉近正样本对的特征距离，推远负样本对的特征距离
\end{itemize}

\begin{figure}[H]
\centering
\includegraphics[width=0.95\textwidth]{slides-images/slide-002.jpg}
\caption{对比学习流程：对同一图片的不同增强应拥有相似的特征向量（绿色），不同图片的增强应远离（红色）\protect\footnotemark}
\end{figure}
\footnotetext{Slides 第3页。}

SimCLR 是对比学习的经典工作之一，但它需要非常大的 batch size 才能收敛。后续的 MoCo（Momentum Contrast）通过维护一个负样本队列和动量编码器解决了这个问题。

\begin{figure}[H]
\centering
\includegraphics[width=0.95\textwidth]{slides-images/slide-004.jpg}
\caption{MoCo 通过维护一个负样本队列和动量编码器（EMA 更新），避免了对超大 batch size 的需求\protect\footnotemark}
\end{figure}
\footnotetext{Slides 第5页。}

\begin{knowledgebox}{DINOv2：当前最强的自监督视觉特征}
DINOv2 在 DINO 的基础上将训练数据扩展到 1.42 亿张图片，是目前实践中最常用的自监督视觉特征模型之一。它结合了对比学习和自蒸馏（self-distillation）思想，在多种下游视觉任务上表现优异。如果需要一个通用的视觉特征 backbone，DINOv2 是首选。
\end{knowledgebox}

\subsection{本章小结}

自监督学习通过 pretext task 在无标签数据上学习特征表示，然后迁移到下游任务。对比学习（SimCLR、MoCo、DINO 系列）是目前最成功的自监督方法之一。DINOv2 通过大规模数据和精心设计的训练策略，成为了当前最强的自监督视觉特征模型。

%% ========== Part 2: 生成模型基础 ==========

